\documentclass[UTF8]{ctexart}
\usepackage[a4paper,margin=2.2cm]{geometry}
\usepackage{hyperref}
\usepackage{array}
\usepackage{longtable}
\usepackage{booktabs}
\usepackage{enumitem}
\usepackage{amsmath}

\title{实验报告：构建 Karpathy 风格长上下文知识库}
\author{马汝坡\quad 杭州电子科技大学\quad 大二电子信息工程\quad 学号：24200516}
\date{2026 年 4 月}

\begin{document}
\maketitle

\section{1. 架构对比分析（重点）}
\subsection{1.1 传统 RAG 与 Karpathy 方案对比}
\begin{longtable}{p{0.2\textwidth}p{0.35\textwidth}p{0.35\textwidth}}
\toprule
维度 & 传统 RAG & Karpathy 长上下文方案（本项目） \\
\midrule
数据流程 & 分块、向量化、检索 Top-K、生成 & raw 资料输入、LLM 整库阅读、结构化 Markdown 输出 \\
召回机制 & 依赖检索结果，可能漏召回 & 整库进入上下文，召回更完整 \\
跨文档推理 & 受检索片段限制 & 单次上下文内可直接跨页综合 \\
系统复杂度 & 需要向量库和检索组件 & 文件系统 + LLM API 即可 \\
适用规模 & 中大规模知识库 & 个人/课程级中小规模知识库 \\
\bottomrule
\end{longtable}

\subsection{1.2 本项目设计决策}
\begin{enumerate}[leftmargin=1.5em]
\item 使用 Markdown 作为唯一存储：便于 Obsidian 展示与人工维护。
\item 不引入向量数据库：遵循课程“极简架构”要求。
\item 同时保留两条编译路径：\texttt{build\_kb.py} 用于稳定样例，\texttt{compile\_raw\_llm.py} 用于自动编译原始资料。
\item 问答采用整库上下文输入：保持“非检索、长上下文”原则。
\end{enumerate}

\subsection{1.3 长上下文方案局限性}
\begin{itemize}[leftmargin=1.5em]
\item Token 成本随知识库规模增加而明显上升。
\item 输入变长导致响应时延上升。
\item 当知识库规模超过上下文窗口时，需要分层编译与分主题组织。
\end{itemize}

\section{2. Prompt 工程细节}
\subsection{2.1 Compilation Prompt 结构}
在 \texttt{src/compile\_raw\_llm.py} 中，Prompt 规定模型输出严格 JSON，核心字段为：
\texttt{site\_title, pages[].title, summary, highlights, equations, questions, tags, related}。

关键约束：
\begin{enumerate}[leftmargin=1.5em]
\item pages 数量至少 10。
\item related 必须引用 pages 中存在的标题。
\item 输出语言为中文，主题聚焦电磁场与电磁波课程。
\end{enumerate}

\subsection{2.2 Prompt 迭代优化}
\begin{itemize}[leftmargin=1.5em]
\item 初版问题：输出混入解释文本，JSON 解析失败。\newline 改进：增加“只输出 JSON”并启用 JSON 模式。
\item 第二版问题：页面数量不足。\newline 改进：硬约束 pages \textgreater= 10。
\item 第三版问题：related 指向不存在页面。\newline 改进：后处理校验并过滤无效引用。
\end{itemize}

\section{3. 测试案例}
\subsection{案例一：静电场与高斯定律关系}
问题：高斯定律和静电场有什么关系？\newline
跨文档页面：矢量场、坐标系与库仑定律；高斯定律与电通量；麦克斯韦方程组与统一图景。

\subsection{案例二：位移电流的作用}
问题：位移电流为什么能把静电理论扩展到电磁波理论？\newline
跨文档页面：时变电磁场与电磁感应；麦克斯韦方程组与统一图景；均匀平面波与传播特性。

\subsection{案例三：传输线与平面波联系}
问题：传输线分析和均匀平面波分析在工程上有什么联系？\newline
跨文档页面：均匀平面波与传播特性；传输线、驻波与波导；辐射、天线与远场。

\subsection{编译前后截图说明}
提交版 PDF 需包含：raw 目录截图、wiki 目录截图、问答终端截图、Obsidian 图谱截图。

\section{4. 性能评估}
\subsection{4.1 知识库规模（实测）}
\begin{itemize}[leftmargin=1.5em]
\item Markdown 文件数：11（索引 1 + 章节 10）
\item 总字符数：6297
\item 中文字符数：3018
\item Token 估算：约 3704（按 chars/1.7 粗估）
\end{itemize}

\subsection{4.2 编译耗时与问答响应时间（实测）}
\begin{itemize}[leftmargin=1.5em]
\item \texttt{python src/build\_kb.py}：约 0.095 s
\item \texttt{python src/qa.py "高斯定律和静电场有什么关系"}：约 0.222 s（本地兜底模式）
\end{itemize}

\subsection{4.3 API 调用成本估算}
采用统一估算公式：
\[
\text{cost} = \frac{T_{in}}{10^6} P_{in} + \frac{T_{out}}{10^6} P_{out}
\]
示例：若一次问答输入 4000 token、输出 800 token，按输入 2 元/百万、输出 8 元/百万估算，单次约 0.0144 元。

\section{5. 反思与展望}
\subsection{5.1 适用与不适用场景}
适用：个人学习知识库、课程项目、文档规模可被长上下文覆盖的场景。\newline
不适用：百万字以上知识库、高并发低延迟在线服务。

\subsection{5.2 面向 100 万字规模的改进方向}
\begin{enumerate}[leftmargin=1.5em]
\item 分层编译：先分主题子库，再维护总索引。
\item 增量更新：仅重编译受影响页面。
\item 轻量索引路由：关键词路由到子库后再长上下文回答。
\item 成本控制：批处理编译和结果缓存。
\end{enumerate}

\section{结论}
本项目已完成课程要求中的核心闭环：Markdown 知识库组织、长上下文问答、图谱导出、自动编译入口。后续重点优化方向为：增量编译与真实 API 场景下的系统化评测。

\end{document}
